% 01_restatement.tex —— 问题重述
\section{问题重述}

某工业制造系统需要对 A、B、C、D、E 五个车间开展集中整修任务。每个车间包含若干具有固定先后顺序的工序，工序编号由小到大表示执行顺序。不同工序需要占用不同类型的设备，部分工序需要两类设备共同参与；若一道工序需要两类设备，则两类设备均需完成该工序对应工程量，只有当两类设备均完成后，该工序才可判定完成。

系统中共有两个班组，每个班组配备若干台不同类型设备。设备在不同工序之间可重复使用，但同一时刻每台设备只能服务一道工序。同一设备在同一车间内切换工序时不计运输时间，而跨车间作业时需要计入设备转运时间。转运时间由车间距离和设备移动速度确定，设备移动速度为 $2\,\mathrm{m/s}$，持续作业时间和运输时间均按秒向上取整。

题目要求解决以下四个递进问题：
\begin{enumerate}[label=\textbf{问题\arabic*.}, leftmargin=2.8em]
  \item 仅由班组 1 独立承担 A 车间全部整修任务，求最短完工时长并给出设备调度表。
  \item 仅使用班组 1 设备完成 A--E 五个车间全部任务，求最短完工时长并给出设备调度表。
  \item 使用班组 1 和班组 2 设备共同完成五个车间全部任务，求最短完工时长并给出带班组信息的设备调度表。
  \item 在 50 万元设备购置预算下，综合确定设备购置方案与作业调度策略，使全部任务完成时间最短，并给出调度表和采购表。
\end{enumerate}

该问题本质上是一个带运输时间的多资源调度问题。其难点主要体现在：车间内部工序必须严格按链式顺序执行；同类设备存在多台可替代机器；双设备工序需要同步协调；同一设备跨车间作业时存在序列相关转运时间；最后一问还引入了预算约束下的设备扩容决策。

从调度理论角度看，该问题兼具柔性作业车间调度和资源受限项目调度特征；同设备跨车间运输时间可视为序列相关准备时间。混合整数规划是小规模作业车间调度中常用的精确建模方法\cite{KuBeck2016}，而带序列相关准备时间的调度生成与排序约束也已有较成熟研究\cite{Artigues2005}。本文在完整 MILP 调度模型基础上，引入 CP-SAT 进行关键下界和采购模型验证，主要参考了近年来约束规划在调度问题中的应用思路\cite{PyJobShop2025,ORToolsCPSAT} 以及 SciPy/HiGHS 对 MILP 的求解接口\cite{SciPyMilp}。
